% ==========================================
% فایل: steps/step00.tex
% گام ۰: بارگذاری و پیش‌پردازش اولیه مجموعه‌ها
% ==========================================

\section{گام (۰): بارگذاری و آماده‌سازی تصاویر ورودی}

\begin{stepbox}
\field{هدف و روند کار}{
در این تمرین به جای اجرای دستی کدها برای هر مجموعه تصویر (\lr{Set 1, 2, 3})، یک تابع جامع پایتون توسعه دادیم تا فرآیند خواندن، پردازش و ذخیره‌سازی را خودکار کند. در اولین قدم، تصاویر ورودی خوانده شده و برای نمایش صحیح و همچنین پردازش‌های بعدی، فضای رنگی آن‌ها اصلاح می‌شود.
}

\field{تحلیل کد: پیش‌پردازش فضای رنگی}{
کتابخانه \lr{OpenCV} به صورت پیش‌فرض تصاویر را به جای استاندارد \lr{RGB}، در فضای \lr{BGR} می‌خواند. اگر این فضا اصلاح نشود، در زمان نمایش تصاویر با \lr{Matplotlib} دچار اعوجاج رنگی می‌شویم. کدهای کلیدی این بخش و منطق آن‌ها به شرح زیر است:
}

\begin{codebox}[بارگذاری و تغییر فضای رنگی]
\begin{lstlisting}[language=Python]
# Read image from path (default is BGR format)
img1_bgr = cv2.imread(img1_path)

# Convert BGR to RGB for correct plotting with matplotlib
img1_rgb = cv2.cvtColor(img1_bgr, cv2.COLOR_BGR2RGB)

# Convert to grayscale for feature extraction algorithms
img1_gray = cv2.cvtColor(img1_bgr, cv2.COLOR_BGR2GRAY)
\end{lstlisting}
\end{codebox}
\end{stepbox}